\section{Widths, Scaling Laws, and Mean Lifetimes}
\label{sec:particle_decay}
\subsection{Internal derivation status and obligations}
A dual substrate account should obtain the decay law from the substrate definition and earlier established dynamics. A familiar external weak-interaction calculation is a comparison target, not a substitute for that derivation.

The earlier formalization establishes the lattice geometry, scheduling and selected bookkeeping quantities. It does not yet construct an $E_8$ transition operator whose continuum limit has a canonically normalized fermion kinetic term, a chiral current, a massive mediator and the relevant decay measure. Those steps must be proved before the fifth power can be claimed as an unconditional substrate result.

There is a precise conditional bridge. With the earlier four-dimensional baseline, suppose the internal action yields a field scaling dimension $d$, a four-field coupling dimension $c$, and a width exponent $k$ satisfying
\begin{equation}
2d+1=4,\qquad c+4d=4,\qquad 2c+k=1.
\end{equation}
Then $d=3/2$, $c=-2$ and $k=5$. Lean checks this implication directly using the earlier baseline. The equations about kinetic, interaction and rate scaling remain explicit proof obligations, not imported facts or new global axioms. Each now has a derived source, stated as a scaling degree of an actual formula (\texttt{ScalingObligations}). The kinetic operator of Paper~II's Hamiltonian is linear in momentum and squares to $|k|^2$, so it is first order with unit speed (\texttt{momentum\_linear}, \texttt{dispersion}). $G_F=1/(\sqrt2v^2)$, derived from the electroweak reference of Paper~III, has degree $-2$ (\texttt{fermi\_scaling}). The leptonic width $G_F^2m^5/(192\pi^3)$ derived below has degree $5$ in the mass, and rescaling $G_F$ by $r^{-2}$ and $m$ by $r$ multiplies it by $r$ (\texttt{rate\_scaling}, \texttt{rate\_dimension\_one}), so the three relations hold at $(3/2,-2,5)$ (\texttt{obligations\_met}). What remains assumed is the dimension count of the action, which turns derivative order one into $2d+1=4$, and the contact-interaction inputs of the width.

The lattice geometry alone does not choose the exponent. Lean gives an internal witness: on the explicit $E_8$ lattice, let $m(z)=\langle z,z\rangle/2$. Both $m(z)^3$ and $m(z)^5$ are positive for nonzero $z$, invariant under every norm-preserving symmetry, and equal to one on a chosen root. On twice that root they give $64$ and $1024$. Thus symmetry, positivity and one-root calibration are insufficient; an actual transition law must distinguish them.

\paragraph{Comparison boundary.} Appendix~\ref{app:weak_comparison} records the conventional weak-interaction calculation as a target for the proposed duality. It is not used to establish an $E_8$ transition operator. Until the internal scaling relations above are derived, the reference-width law below remains a stated benchmark model.

\subsection{A reference-width model}
Retain the illustrative ansatz
\begin{equation}\label{eq:reference_width}
\Gamma(M)=\Gamma_0\left(\frac{M}{M_0}\right)^5|V|^2CF,
\end{equation}
with equal-channel multiplicity $C$ and phase-space factor $0<F\leq1$. At fixed other inputs, the lifetime scales inversely with each width factor. In particular, reducing $F$ increases the lifetime: it cannot correct an estimate that is already too long while all other inputs remain fixed.

The benchmark multiplicities $5$ and $9$ do not include a derivation of hadronic amplitudes or radiative corrections. Table~\ref{tab:lifetimes} retains the original rounded numbers, with their nonuniform accuracy exposed. These are mean lifetimes, not half-lives.
\begin{table*}[t]
\centering
\caption{Illustrative draft benchmarks. Heavy-quark-sector comparisons require a specified hadron species and mass convention before precision testing. References are the rounded draft inputs, not a new experimental fit.}
\label{tab:lifetimes}
\begin{tabular}{lrrr}\toprule
State or sector & Model $\tau$ (s) & Draft reference $\tau$ (s) & Model/reference \\\midrule
Muon & $2.20\times10^{-6}$ & $2.19\times10^{-6}$ & $1.005$ \\
Tau & $3.27\times10^{-13}$ & $2.90\times10^{-13}$ & $1.128$ \\
$b$ sector & $1.49\times10^{-12}$ & $1.53\times10^{-12}$ & $0.974$ \\
$c$ sector & $1.84\times10^{-12}$ & $1.05\times10^{-12}$ & $1.752$ \\
$K^+$ & $3.25\times10^{-8}$ & $1.24\times10^{-8}$ & $2.621$ \\
Neutron & $8.42\times10^2$ & $8.77\times10^2$ & $0.960$ \\\bottomrule
\end{tabular}
\end{table*}

\paragraph{Lepton universality check.} For the leptonic modes the coefficient of the fifth power is not a free benchmark. Given the spin-averaged V-A contact amplitude, the covariant form of the neutrino-pair integral and the massless phase-space measure, $\Gamma(\ell\to e\nu\bar\nu)=G_F^2m_\ell^5/(192\pi^3)$ with the spectrum $2x^2(3-2x)$; Lean checks each step of this reduction (\texttt{MuonRate.muon\_rate}, \texttt{spectrum\_form}). The amplitude itself follows from the trace algebra of the Dirac matrices reconstructed from Paper~II's Hamiltonian: $\sum_{\mu\nu}\operatorname{Tr}[\not a\gamma^\mu\not b\gamma^\nu(1-\gamma^5)]\operatorname{Tr}[\not c\gamma_\mu\not d\gamma_\nu(1-\gamma^5)]=64(a\cdot c)(b\cdot d)$, which gives $64G_F^2(p\cdot\bar k)(p_e\cdot k)$ after the muon spin average (Paper~II, \texttt{VATrace.muon\_amplitude}; \texttt{MuonAmplitude.amplitude\_matches}). For the spin sums, $\not p\gamma^0$ is Hermitian, satisfies $(\not p\gamma^0)^2=2p^0\,\not p\gamma^0$ for $p^2=0$, is annihilated by $\not p$ and has trace $4p^0$, so it is $2p^0$ times the rank-two projector onto the massless solutions (\texttt{MuonAmplitude.spin\_projector}, \texttt{spin\_solutions}, \texttt{spin\_hermitian}, \texttt{spin\_trace}); identifying it with $\sum_s u_su_s^\dagger$ for any orthonormal basis of the solutions, massive or massless, is proved in Paper~II's research module \texttt{SpinorCompleteness} ($\sum u\bar u=\not p+m$, $\sum v\bar v=\not p-m$). The muon is massive, but its mass term drops out of the V--A trace exactly, because $(1-\gamma^5)\gamma^\nu(1-\gamma^5)=0$ (\texttt{MassiveMuon.chiral\_sandwich}, \texttt{massive\_amplitude}); $m_\mu$ enters the rate only through the kinematics. With one $G_F$ for every lepton (the shared electroweak reference of Paper~III), the ratio of the tau and muon leptonic widths cancels $G_F$ and $192\pi^3$ (\texttt{TauLifetime.rate\_ratio}), so $\tau_\tau=\tau_\mu(m_\mu/m_\tau)^5B(\tau\to e\nu\bar\nu)$. With the derived lepton masses, the measured $\tau_\mu$ and $B_e=17.82\%$, this gives $291.15$--$291.17$~fs at tree level (\texttt{tau\_lifetime\_tree}); the standard small corrections, entered as a factor in $[0.9995,0.9997]$, place it $0.8$--$1.0\sigma$ above the measured $290.3\pm0.5$~fs (\texttt{tau\_lifetime\_pull}). This replaces the tau benchmark row of Table~\ref{tab:lifetimes} for the leptonic channel. It does not decide the $2\sigma$ tension of the tau mass formula: that tension moves $\tau_\tau$ by only about $0.15$~fs.

\subsection{Interaction classes and proposed routing gears}
Strong, electromagnetic and weak processes provide qualitative scales, not four exact disjoint lifetime bands. Widths depend on amplitudes, thresholds, masses and exclusive channel sums. A power $M^n$ is a conditional scaling law, not a universal interface classifier. Prompt $W$, $Z$ and top decays should not be classified as strong merely because their lifetimes are short. The proposed routing gears remain organizing hypotheses to be tested against a specified interaction law.

\subsection{Weak-boson and environmental ratios}
The geometric assignment $\sigma/(\sigma+1)$ gives $5/6$ at the baseline. Identifying this count with $\Gamma_W/\Gamma_Z$ is a separate phenomenological hypothesis; it does not calculate electroweak partial widths.

For the draft's neutron environment model, retain
\begin{equation}
\mathcal D=\frac{\Delta}{\Delta-\chi/\sigma}
=\frac{43}{43-2/5}=\frac{215}{213}>1.
\end{equation}
If $\Gamma_{\rm beam}=\Gamma_{\rm bottle}/\mathcal D$ is stipulated, then $\tau_{\rm beam}=\mathcal D\tau_{\rm bottle}$. Lean checks this reciprocal relation. The draft input $877.75$ s gives $885.9918$ s. This rescaling is not a demonstration that environmental width changes cause the beam--bottle discrepancy; that claim needs a microscopic mechanism, controls and an independent uncertainty analysis.

\subsection{Formal scope}
The Lean project checks conserved-label routing, finite latency descent, survival identities under a constant hazard, partial-width addition, branching normalization, dimensional and reference-width scaling, and the environmental factor. It does not derive amplitudes, a particle spectrum, confinement, an exact four-gear classification or a quantum Yang--Mills mass gap from $E_8$.
